\documentclass[10pt,aspectratio=169]{beamer}
\input{../../_en_preambulo_beamer}
% Comandos y macros estadísticas compartidas para las presentaciones Beamer en inglés (Metropolis)

% Conjuntos numéricos básicos
\newcommand{\R}{\mathbb{R}}
\newcommand{\N}{\mathbb{N}}
\newcommand{\Q}{\mathbb{Q}}
\newcommand{\Z}{\mathbb{Z}}
\newcommand{\set}[1]{\left\{ #1 \right\}}
\newcommand{\abs}[1]{\left|#1\right|}
\newcommand{\norm}[1]{\left\| #1 \right\|}

% Comandos estadísticos y probabilísticos del curso
\newcommand{\Var}{\operatorname{Var}}
\newcommand{\cov}{\operatorname{Cov}}
\newcommand{\corr}{\rho}
\newcommand{\comb}[2]{\begin{pmatrix} #1 \\ #2 \end{pmatrix}}
\newcommand{\s}{\sigma}
\newcommand{\particion}{\mathfrak{P}}
\newcommand{\card}[1]{\#\left( #1 \right)}
\newcommand{\seq}[3]{\set{#1}_{#2}^{#3}}
\newcommand{\E}{\mathbb{E}}
\newcommand{\Prob}{\mathbb{P}}

% Entornos pedagógicos y cajas en inglés académico natural (soportando tanto nombres en inglés como en español)
\newenvironment{discussion}{%
  \begin{alertblock}{Discussion Question}%
}{%
  \end{alertblock}%
}
\newenvironment{discusion}{%
  \begin{alertblock}{Discussion Question}%
}{%
  \end{alertblock}%
}

\newenvironment{reflection}{%
  \begin{alertblock}{Classroom Reflection}%
}{%
  \end{alertblock}%
}
\newenvironment{reflexion}{%
  \begin{alertblock}{Classroom Reflection}%
}{%
  \end{alertblock}%
}

\newenvironment{paradox}{%
  \begin{alertblock}{Exploratory Paradox}%
}{%
  \end{alertblock}%
}
\newenvironment{paradoja}{%
  \begin{alertblock}{Exploratory Paradox}%
}{%
  \end{alertblock}%
}

\newenvironment{motivation}{%
  \begin{exampleblock}{Motivation: Data Science in Practice}%
}{%
  \end{exampleblock}%
}
\newenvironment{motivacion}{%
  \begin{exampleblock}{Motivation: Data Science in Practice}%
}{%
  \end{exampleblock}%
}

\newenvironment{quickchallenge}{%
  \begin{block}{Quick Team Challenge}%
}{%
  \end{block}%
}
\newenvironment{retoclase}{%
  \begin{block}{Quick Team Challenge}%
}{%
  \end{block}%
}


\title[Transformation of Variables]{Section 5.1: Transformation of Variables}
\subtitle{Distributions of $Y=g(X)$}
\author[J. Castillo Colmenares]{Juliho Castillo Colmenares}
\institute{Tecnológico de Monterrey}
\date{}

\begin{document}
\begin{frame}[plain]\vspace{-0.8cm}\titlepage\end{frame}

\begin{frame}{Roadmap --- Chapter 5}
  \footnotesize
  \begin{itemize}\itemsep=0.08cm
    \item \textbf{05.00} Inferential introduction
    \item \textbf{05.01} Transformation of variables \emph{(today)}
    \item \textbf{05.02} Functions of random variables
    \item \textbf{05.03--05.05} $\chi^2$, $t$, and $F$ distributions
  \end{itemize}
  \begin{block}{Learning objective}
    Obtain the density of $Y=g(X)$ through affine maps, inverse functions, and
    Jacobians, including the non-monotone case.
  \end{block}
\end{frame}

\begin{frame}{Motivation: changing scale}
  \footnotesize
  If $X$ has a known distribution and $Y=g(X)$, we need the distribution of $Y$
  to analyze transformed data.
  \begin{itemize}\itemsep=0.12cm
    \item A transformation may normalize data.
    \item It may stabilize variance or compress extreme values.
    \item It must preserve total probability.
  \end{itemize}
\end{frame}

\begin{frame}{Affine transformation}
  \footnotesize
  If $Y=aX+b$, $a\neq0$, and $X$ has mean $\mu_X$ and variance $\sigma_X^2$,
  \[
    \mu_Y=a\mu_X+b,\qquad \sigma_Y^2=a^2\sigma_X^2.
  \]
  \begin{alertblock}{Normal case}
    If $X\sim N(\mu,\sigma^2)$, then $Y\sim N(a\mu+b,a^2\sigma^2)$.
  \end{alertblock}
\end{frame}

\begin{frame}{Monotone change of variable}
  \footnotesize
  For strictly monotone differentiable $g$ and $Y=g(X)$,
  \[
    f_Y(y)=f_X\!\left(g^{-1}(y)\right)
    \left|\frac{d}{dy}g^{-1}(y)\right|.
  \]
  First write $F_Y(y)=F_X(g^{-1}(y))$, then differentiate by the chain rule.
\end{frame}

\begin{frame}{Inverse Jacobian}
  \footnotesize
  With $x=g^{-1}(y)$ and $J(y)=dx/dy$,
  \[
    f_Y(y)=f_X(x)|J(y)|.
  \]
  \begin{exampleblock}{Geometric reading}
    The absolute Jacobian corrects the local scale change between $x$ and $y$.
  \end{exampleblock}
\end{frame}

\begin{frame}{Non-monotone case}
  \footnotesize
  If several $x$ values produce the same $y$, add every preimage:
  \[
    f_Y(y)=\sum_{x:g(x)=y}f_X(x)\left|\frac{dx}{dy}\right|.
  \]
  \begin{alertblock}{Key difference}
    A single inverse branch is not enough when the function folds over itself.
  \end{alertblock}
\end{frame}

\begin{frame}{Common transformations}
  \footnotesize
  The reading lists transformations useful for positive or skewed data:
  \begin{itemize}\itemsep=0.11cm
    \item $Y=\ln X$ compresses large values.
    \item $Y=\sqrt X$ smooths count-data skewness.
    \item $Y=X^{1/3}$ handles stronger skewness.
    \item Box--Cox: $Y=(X^\lambda-1)/\lambda$.
  \end{itemize}
\end{frame}

\begin{frame}{Exact application 1/4: normal transformation}
  \footnotesize
  For $X\sim N(\mu,\sigma^2)$ and $Y=aX+b$,
  \[
    Y\sim N(a\mu+b,a^2\sigma^2).
  \]
  An affine transformation shifts the mean and scales the variance while
  preserving the normal family.
\end{frame}

\begin{frame}{Exact application 1/4: solution}
  \footnotesize
  \begin{enumerate}\itemsep=0.13cm
    \item identify $a$ and $b$;
    \item transform the mean to $a\mu+b$;
    \item transform the variance to $a^2\sigma^2$.
  \end{enumerate}
  This is the reading's example of normal reproducibility.
\end{frame}

\begin{frame}{Exact application 2/4: exponential to Pareto}
  \footnotesize
  Let $X\sim\operatorname{Exp}(1)$ and $Y=e^X$. Then $g^{-1}(y)=\ln y$ and
  $d(\ln y)/dy=1/y$.
  \[
    f_Y(y)=e^{-\ln y}\frac1y=\frac1{y^2},\qquad y>1.
  \]
\end{frame}

\begin{frame}{Exact application 2/4: solution}
  \footnotesize
  Since $X>0$ and the map is increasing, $Y>1$.
  \begin{alertblock}{Result}
    The density $1/y^2$ for $y>1$ is a Pareto distribution with $\alpha=1$.
  \end{alertblock}
\end{frame}

\begin{frame}{Exact application 3/4: log-normal transformation}
  \footnotesize
  If $X\sim N(\mu,\sigma^2)$ and $Y=e^X$, then $y>0$ and
  \[
    f_Y(y)=\frac{1}{\sigma\sqrt{2\pi}}
    \exp\!\left[-\frac{(\ln y-\mu)^2}{2\sigma^2}\right]\frac1y.
  \]
\end{frame}

\begin{frame}{Exact application 3/4: solution}
  \footnotesize
  \begin{block}{Reading result}
    $Y$ follows a log-normal distribution $LN(\mu,\sigma^2)$.
  \end{block}
  The factor $1/y$ is the inverse Jacobian and must not be omitted.
\end{frame}

\begin{frame}{Exact application 4/4: non-monotone sine}
  \footnotesize
  For $X\sim U(0,2\pi)$ and $Y=\sin X$, each $y\in(-1,1)$ has two preimages:
  \[
    x_1=\arcsin y,\qquad x_2=\pi-\arcsin y.
  \]
  Each branch contributes $1/(2\pi\sqrt{1-y^2})$.
\end{frame}

\begin{frame}{Exact application 4/4: solution}
  \footnotesize
  \[
    f_Y(y)=\frac{1}{\pi\sqrt{1-y^2}},\qquad -1<y<1.
  \]
  \begin{alertblock}{Result}
    This is the arcsine density; dropping one preimage would give only half the
    correct probability.
  \end{alertblock}
\end{frame}

\begin{frame}{Computational bridge 1/4: support}
  \footnotesize
  \begin{center}
    $X>0\xrightarrow{e^x}Y>1$,\qquad
    $X\in(0,2\pi)\xrightarrow{\sin}Y\in(-1,1)$.
  \end{center}
  The support is a quick check for a transformed density.
\end{frame}

\begin{frame}{Computational bridge 2/4: Jacobian}
  \footnotesize
  \[
    J(y)=\frac{d}{dy}g^{-1}(y).
  \]
  \begin{itemize}\itemsep=0.12cm
    \item $g^{-1}(y)=\ln y\Rightarrow |J(y)|=1/y$.
    \item $g^{-1}(y)=\arcsin y\Rightarrow |J(y)|=1/\sqrt{1-y^2}$.
  \end{itemize}
\end{frame}

\begin{frame}{Computational bridge 3/4: conservation}
  \footnotesize
  Each transformed density must integrate to one over its support.
  \begin{block}{Consistency check}
    The Jacobian redistributes mass; it does not create or remove probability.
  \end{block}
\end{frame}

\begin{frame}{Computational bridge 4/4: method choice}
  \footnotesize
  \begin{tabular}{ll}
    \hline
    Shape of $g$ & Action \\
    \hline
    Affine & transform mean and variance \\
    Monotone & one inverse and one Jacobian \\
    Non-monotone & add all branches \\
    \hline
  \end{tabular}
  \begin{alertblock}{Practical rule}
    The function's geometry determines the solution's form.
  \end{alertblock}
\end{frame}

\begin{frame}{Synthesis of the section}
  \footnotesize
  \begin{itemize}\itemsep=0.12cm
    \item Affine transformations preserve normality.
    \item Change of variable uses an inverse and its Jacobian.
    \item Non-monotone maps require all preimages.
    \item Support and total mass verify the result.
  \end{itemize}
\end{frame}

\begin{frame}{Curricular transition}
  \footnotesize
  \begin{block}{Established}
    We can derive distributions for transformations of one random variable.
  \end{block}
  \begin{alertblock}{Next section}
    \textbf{05.02 Probability distributions of functions of a random variable}:
    extend the method to one or several variables and level surfaces.
  \end{alertblock}
\end{frame}
\end{document}
